\section{Rank-seven neighborhoods of rank-six spectra}

\begin{theorem}\label{thm:rank6-neighborhood}
For each \(g>0\) there is an explicit \(\varepsilon(g)>0\),
independent of the qubit bipartition, such that the sorted
computational representative is globally optimal at order two
for positive ordered rank-seven probabilities with \(t_0=1\),
\(\min(t_0-t_1,t_1-t_2)\ge g\), and \(t_6\le\varepsilon(g)\).
With strict inequality \(t_6<\varepsilon(g)\) and distinct
positive probabilities, the optimizers are precisely the local
Clifford images of the sorted representative up to phase.
\end{theorem}

Let \(C\) be positive with rank at most \(R\), maximum eigenvalue
one, and \(c_W=\tr(CW)\). Write \(\eta=\Delta_d(C)/d\) and put
\(\rho=1/(4R^4)\). By the defect formula and
\(\tr(C^4)\le[\tr(C^2)]^2\) one has
\[
 d^{-1}\sum_W c_W^4\ge [\tr(C^2)]^2-\eta.
\]
Also \(\sum_W c_W^2=d\tr(C^2)\) and \(|c_W|\le R\). At
\(\eta\le1/4\), the set \(S=\{W:c_W^2>1/2\}\) therefore has
at least \(\rho d\) labels. The anticommuting defect bound
of \cite[Theorem~14]{sierant} gives at most \(2\eta d^2\)
ordered anticommuting pairs in \(S\).

Choose \(k=\lceil24/\rho^2\rceil+1\) and
\(\zeta=1/[16(k+1)]\). At \(\eta\le\zeta\rho^2/4\), at
least half of \(S\) consists of labels that each anticommute
with at most \(\zeta|S|\) labels of \(S\). Call this subset
\(T\), which contains the identity label zero. Every binary
sum in \((k+1)T\) has at most \((k+1)\zeta|S|\) anticommuting
partners in \(S\). Its unnormalized symplectic Fourier
coefficient of \(1_S\) is at least \(|S|/2\). Parseval on the
\(d^2\) labels gives
\[
 |(k+1)T|\le4d^2/|S|\le4d/\rho .
\]
If a shortest word in \(T\) for an element of \(\langle T\rangle\)
has length at least \(k\), the translates of \(T\) by every
third prefix of its first \(k\) steps are disjoint. An
intersection would shorten the word by replacing at least
three steps by two. These translates lie in \((k+1)T\), giving
more than \(4d/\rho\) elements by the choice of \(k\), a
contradiction. Thus \(V=\langle T\rangle\subseteq kT\).
Every label of \(V\) anticommutes with at most \(|S|/16\)
labels of \(S\), so averaging over \(V\) shows
\(|S\setminus V^\perp|\le|S|/8\).
Consequently \(B=T\cap V^\perp\) is commuting and has
\(|B|\ge3\rho d/8\).

Let \(H=\langle B\rangle\),
\(m=\lceil\log_2(8/\rho)\rceil\), and \(\beta=\rho/[16(m+1)]\).
Call a character of \(H\) sparse if it is negative on fewer
than \(\beta d\) elements of \(B\). There cannot be \(m\)
independent sparse characters, since their common kernel
would have cardinality at most \(d/2^m\le\rho d/8\) and
contain more than \(\rho d/8\) elements of \(B\).
Taking the common kernel \(H'\) of all sparse characters
gives joint eigenspaces of dimension
\(\ell=d/|H'|\le L=4\cdot2^m/\rho\). Each Pauli outside
\((H')^\perp\) anticommutes with at least \(\beta d\) labels
of \(B\). Those labels have \(c_W^2>1/2\). The defect bound
and Parseval then give
\[
 \|C-\mathcal E_{H'}C\|_{\rm HS}^2\le A\eta,\qquad
 A=1/(2\beta),
\]
where \(\mathcal E_{H'}\) is Pauli conjugation averaging.

For the logical Pauli labels \(p\in (H')^\perp/H'\), define
\(a_p^2=|H'|^{-1}\sum_{W\in p+H'}c_W^2\).
Parseval over the joint characters of \(H'\) identifies this
with the sum over blocks of the squared logical coefficients
at \(p\). If two logical labels anticommute, their entire
cosets anticommute, giving \(\eta\ge4a_p^2a_q^2/\ell^2\).
For \(s=g/(16L)\) and \(\eta<2s^4/L^2\), the labels with
\(a_p>s\) therefore commute. Let \(G\) be the group generated
by their cosets together with \(H'\). Retaining just these high
cosets in \(\mathcal E_{H'}C\) gives a matrix \(C_{\rm hi}\)
that is scalar on the joint eigenspaces of \(G\). The omitted
terms change the operator norm in each \(H'\) block by at
most \(\ell s\le g/16\), because every omitted coefficient
in the block has modulus at most \(s\). If
\(\eta<g^2/(16A)\), then
\(\|C-C_{\rm hi}\|<g/4+g/16<g/2\).

Each label in \((H')^\perp\setminus G^\perp\)
anticommutes with some entire high coset. Counting these
pairs by the endpoint outside \(G^\perp\) and using the defect bound
gives
\[
 \|C-\mathcal E_G C\|_{\rm HS}^2
 \le \mathcal A\eta,\qquad
 \mathcal A=A+L/(4s^2).
\]
The projections in this bound are nested and the two sets of
removed Pauli coefficients are disjoint.

With \(t_0-t_1\ge g\), the bound on \(C-C_{\rm hi}\) makes the
largest eigenvalue of \(C_{\rm hi}\) simple. All joint
eigenspaces of \(G\) have equal dimension, so \(G\) is maximal.
The gaps of both leading eigenlines \(P_0,P_1\) of \(C\) are
at least \(g\). Comparing commutators in its eigenbasis gives
\(\|P_j-\mathcal E_GP_j\|_{\rm HS}^2\le\mathcal A\eta/g^2\).
For a rank-one projector, its largest basis weight in the
stabilizer basis of \(G\) is at least one minus this
dephasing loss. At \(\eta<g^2/(4\mathcal A)\) the two chosen
reference lines must be distinct by orthogonality. Thus,
writing
\[
 \eta_*=\min\{1/4,\zeta\rho^2/4,g^2/(16A),2s^4/L^2,
                           g^2/(4\mathcal A)\},
 \qquad
 K(g,R)=2/\eta_*+2\mathcal A/g^2,
\]
the sum of the squared infidelities is at most \(K(g,R)\eta\).
This holds for \(\eta\ge\eta_*\) as well because the sum is
always at most two.

\begin{proof}[Proof of Theorem~\ref{thm:rank6-neighborhood}]
In the rank-seven forward comparison
\eqref{eq:forward-dominant}, \(t_0=1\) gives
\(0\le\gamma_{02}\le\gamma_{01}\le t_6+t_6^2\le2t_6\).
Lemma~\ref{lem:compensated-anchors} and the ordered-line
rounding give the upper bound
\[
  \bigl(336\cdot126500\,t_6 K(g,7)-1\bigr)\eta
\]
on the purity excess. In particular
\(\varepsilon(g)=[336\cdot126500\,K(g,7)]^{-1}\) suffices.
Apply the two-frame comparison and stabilizer padding.
For a strict cutoff the defect must vanish at any optimal
one-frame objective. At an optimizing physical state both
frames attain equality, so each physical marginal has
commuting Pauli support. The rigidity statement of
Theorem~\ref{thm:one-stabilizer-marginal} gives the equality
classification for distinct probabilities.
\end{proof}

\begin{theorem}\label{thm:doublet-boundary}
For each \(g>0\) there is a dimension-independent positive
cutoff \(\varepsilon_{\rm eq}(g)\) such that sorted rank-seven
probabilities with \(t_0=t_1=1\), \(1-t_2\ge g\), and
\(t_6\le\varepsilon_{\rm eq}(g)\) have the same global optimum.
\end{theorem}

\begin{proof}
Use the same constants and proceed to \(\mathcal E_G C\) as above
with \(\eta<\eta_*\). Put \(z=(t_1+t_2)/2\). Since
\(\|C-C_{\rm hi}\|<5g/16\), the spectrum of \(C_{\rm hi}\) has
exactly two values above \(z\), counting multiplicities, and its
distance from \(z\) is at least \(3g/16\). Its joint \(G\) block
size must divide two. Hence its two-dimensional upper spectral
space \(Q\) is either a single rank-two stabilizer block or a
span of two members of a common stabilizer basis. In both cases
it is a rank-two stabilizer code.

Put \(D=\mathcal E_G C\). The difference \(D-C_{\rm hi}\) has
operator norm at most \(g/16\), by averaging
\(\mathcal E_{H'}C-C_{\rm hi}\) over \(G\). The operators are
block diagonal and \(C_{\rm hi}\) is scalar on each \(G\) block,
so \(Q\) is also the upper spectral subspace of \(D\).
For an eigenvector \(x_i\) of \(C\) in its leading plane \(P\)
and an eigenvector \(u_j\) of \(D\) below \(z\), with eigenvalue
\(\lambda_j\), the identity
\((t_i-\lambda_j)\langle u_j|x_i\rangle
=\langle u_j|(C-D)x_i\rangle\) gives
\[
 1-\tfrac12\tr(PQ)\le 2\|C-D\|_{\rm HS}^2/g^2
                    \le 2\mathcal A\eta/g^2.
\]
The rank-two stability estimate \eqref{eq:delta-dimensionfree}
gives \((L_{01}-2)_+\le20\mathcal A\eta/g^2\).
Here \(b_0=0\), so the forward coefficient \(\gamma_{02}\)
vanishes and \(\gamma_{01}\le t_6^2\).
In the near case the purity excess is therefore nonpositive
if \(t_6\le g/\sqrt{3360\mathcal A}\). In the far case
\(\eta\ge\eta_*\), use \(L_{01}\le5/2\), which suffices for
\(t_6\le\sqrt{\eta_*/84}\). The smaller of the two cutoffs works
on every bipartition by the paired comparison.
\end{proof}

\begin{theorem}\label{thm:simple-top-boundary}
For each \(g>0\) there is a dimension-independent positive
cutoff \(\varepsilon_{\rm top}(g)\) such that sorted rank-seven
probabilities with \(t_0=1\), \(t_0-t_1\ge g\), and
\(t_6\le\varepsilon_{\rm top}(g)\) satisfy global optimality.
The remaining probabilities need no separation.
\end{theorem}

\begin{proof}
For \(\eta<\eta_*\), the preceding construction gives a maximal
\(G\) using just the leading gap. Let \(v_0\) be the stabilizer
reference of \(x_0\), with squared infidelity at most
\(e_0=\mathcal A\eta/g^2<1/4\).

For a two-wedge with reference stabilizer lines \(v,w\) and
squared infidelities \(e_x,e_y\), retain the two one-anchor
self-expectations in the proof of
\eqref{eq:delta-dimensionfree}. Specifically, let \(a^2,b_0^2,b_1^2\)
be the self-expectations of the component on \(v\wedge w\) and the
two components \(\eta_0,\eta_1\) containing one reference line.
Using the star compressions of the operator \(L=2I-M_Y\) gives
\(a^2+b_0^2+b_1^2-2a(b_0+b_1)\ge-2b_0b_1\).
Since \(L\preceq2I\), this is at least
\(-4\|\eta_0\|\|\eta_1\|\). The other cross terms use
\(\|L-3I/4\|\le5/4\), and the component \(z\) outside both
reference lines has self-expectation at least \(-\|z\|^2/2\).
With \(\|\eta_0\|\le2\sqrt{e_y}\),
\(\|\eta_1\|\le2\sqrt{e_x}\), and \(\|z\|\le\sqrt{e_x e_y}\),
the result is \((L_{xy}-2)_+\le 29\sqrt{e_x e_y}\).

Write the forward bracket in \eqref{eq:forward-dominant} as
\[
 \gamma_{02}(L_{01}+L_{02}-4)
   +t_6^2 b_1(t_0-t_3)(L_{01}-2).
\]
For \(b_1>0\), the line \(x_1\) has gap at least
\(h_1=\min(g,b_1)\). Its dephasing loss is at most
\(d_1=\mathcal A\eta/h_1^2\).
If \(d_1<1/2\), its most likely reference line is distinct
from \(v_0\). Otherwise any distinct reference has
infidelity at most \(1\le2d_1\). The asymmetric estimate and
\(b_1/h_1\le1/g\) therefore bound the second term by
\(29\sqrt2\,t_6^2\mathcal A\eta/g^2\). The term vanishes if \(b_1=0\).

If \(b_2>0\), set \(P=\operatorname{span}\{x_1,x_2\}\).
Its gap to the complementary eigenvalues is at least
\(h_2=\min(g,b_2)\), so its projector's squared dephasing
loss is at most \(d_2=\mathcal A\eta/h_2^2\).
For \(d_2<1/4\) the top two coordinate masses of its
projector sum to at least \(2-d_2\). Both coordinates are
different from \(v_0\), whose mass is at most \(e_0\). Their
span \(Q\) admits the polar isometry \(PQ(QPQ)^{-1/2}\)
onto \(P\). It gives an orthonormal basis of \(P\) with
squared infidelities at most \(2d_2\).
If \(d_2\ge1/4\), any two distinct references outside \(v_0\)
and any orthonormal basis of \(P\) give errors at most
\(1\le4d_2\).
The sum \(L_{01}+L_{02}\) is invariant under this basis
change inside \(P\), since each overlap trace counts a
quadruple with the number of legs in \(P\). Thus the
asymmetric estimate on the two legs gives
\(L_{01}+L_{02}-4\le116\mathcal A\eta/(g h_2)\).
As \(\gamma_{02}\le b_2t_6\) and \(b_2/h_2\le1/g\),
the first term costs at most \(116t_6\mathcal A\eta/g^2\).
It vanishes if \(b_2=0\).

The combined bracket is at most \(160t_6\mathcal A\eta/g^2\).
For \(\eta\ge\eta_*\) use \(L_{ij}\le5/2\) to bound it instead
by \(2t_6\). Paying the coefficient \(168\) with the defect
\(\eta\) shows that
\[
 \varepsilon_{\rm top}(g)=\min\{\eta_*/336,\ g^2/(26880\mathcal A)\}
\]
suffices. The physical comparison and equality argument for
a strict cutoff are the same as in Theorem~\ref{thm:rank6-neighborhood}.
\end{proof}

\begin{lemma}\label{lem:affine-eight-anchor}
A fixed two-plane contained in the coordinate space of a binary
affine three-space has normalized five-set occupancy at most two
against every orthogonal five-space.
\end{lemma}

\begin{proof}
Let \(H\) be the affine three-space and \(X\) the fixed two-plane.
For a direction \(L\), twice contracting its stabilizer four-form
by \(X\) yields a two-form in a single block of \(X^\perp\).
The blocks are the six-dimensional space \(\mathbb C^H\ominus X\)
and the eight-dimensional coordinate spaces of the other cosets
of \(\operatorname{dir}H\). Indeed a contributing plane meets
\(H\) in either four points, or two points with the other two in
a coset fixed by \(L\). Thus the two-form operator is an orthogonal
sum of operators on the two-exteriors of these blocks.
In an orthonormal basis of a block let \(p_{ij}\) be its nonnegative
coordinate pair weights. The one-hole isometry and the five-plane
bound of \cite[Theorem~17 and Corollary~4]{sierant} give degree
one at each vertex and triangle sums at most one.
For a subset \(A\) in a block of even size \(m=6\) or \(8\),
the total internal weight \(e(A)\) satisfies
\(e(A)-e(A^c)=|A|-m/2\) by the degree identity.
For even sizes the degree bound gives \(e(A)\le |A|/2\).
For odd sizes, singleton and triangle bounds and their complements
give \(e(A)\le\lfloor |A|/2\rfloor\).

Choose the block bases to diagonalize the compressions of the
five-plane projection to the blocks, with diagonal values \(s_i\).
The two-exterior compression then gives occupancy
\(\sum_b\sum_{i<j\in b}p_{ij}s_i s_j\).
In the union basis, the squared coefficients of a unit volume form
of the five-plane give a distribution on five-element subsets
\(S\). The inclusion probability for two labels in the same block
is the corresponding two-by-two principal minor of the projection,
which is \(s_i s_j\). The occupancy is therefore the expectation
of \(\sum_b e(S\cap b)\), bounded by two for every five-set.
\end{proof}

\begin{theorem}\label{thm:all-rank6-boundary}
For each \(s\in(0,1]\) there is an explicit dimension-independent
\(r_6(s)>0\) such that normalized rank-seven probabilities with
\(t_5\ge s\) and \(t_6\le r_6(s)\) have sorted computational-basis
optimality on every qubit bipartition and over all local unitaries.
In particular every rank-six spectrum has a relative rank-seven
neighborhood of global optimality.
\end{theorem}

\begin{proof}
Use the high-coset construction with \(g=s/2\), and put
\(\eta_0=\min\{\eta_*,g^2s^8/(20\cdot560^2\mathcal A)\}\).
Require \(t_6\le s/2\). For \(\eta<\eta_0\), use the truncation at
\(z=(t_5+t_6)/2\). Exactly six of its eigenvalues are above \(z\),
at distance at least \(3g/16\) from the cutoff. The common \(G\)
block size is a power of two dividing six, hence one or two.
The same upper spectral projection \(Q\) applies to
\(D=\mathcal E_G C\), because \(D\) preserves the blocks and its
distance to \(C_{\rm hi}\) is at most \(g/16\).

After a Clifford, \(Q\) consists of six computational labels.
For block size two these are three coordinate cosets of one
binary direction. The three points in the quotient lie in one
affine plane, whose preimage is an eight-point affine space.
For block size one the six labels can span a larger affine space.
Let \(P\) be the actual leading six-space.
The spectral-projection residual as in
Theorem~\ref{thm:doublet-boundary} gives
\(\tr P(I-Q)\le4\mathcal A\eta/g^2\).
The direct principal-angle rotation \(e^K\) maps \(Q\) to \(P\)
with \(K\) off-diagonal in \(Q\oplus Q^\perp\) and
\(h^2=\|K\|_{\rm op}^2\le \pi^2\mathcal A\eta/g^2\).
Pull back the actual frame by \(e^{-K}\), preserving its
eigenvalues. If the six reference labels lie in no affine
three-space, they contain at most one affine plane: two such
planes inside six labels would meet in a line and span an affine
three-space. Hence the compression of \(\Pi_{\rm F}\) to their
four-exterior space is zero or \(q\) times a rank-one projector.
The reference \(c_{\rm F}/q\) is at most \(t_0t_1t_2t_3+20t_6\),
using the one-hole bound for the twenty tail wedges. The sorted
reference has the planes \(0123,0145,2345\), giving at least
\(t_0t_1t_2t_3+2s^4\). Transporting the frame changes
\(c_{\rm F}/q\) by at most \(280h\): there are 35 terms and each
normalized squared projection changes at speed at most \(8h\).
At \(\eta<\eta_0\) and \(t_6\le s^4/40\) the two errors are each
at most \(s^4/2\). This gives strict exterior comparison, and
therefore strict purity comparison by the rank-seven compression.

It remains that \(Q\) lies in an affine eight-point space \(H\).
Translate the labels so that \(H\) is linear.
Let \(a\) be the difference of its two missing labels, so \(Q\)
is three \(a\)-cosets of \(H\) and the remaining coset is \(E\).
The two fixed pairs \(01,02\) are in \(H\), so each reference
occupancy is at most two by Lemma~\ref{lem:affine-eight-anchor}.

Every projected four-wedge has zero derivative of its squared
norm along this rotation at the reference. To see it, compare
the orthogonal direction labels of affine planes contributing
to its value and to its derivative. A wedge of four leading
vectors uses directions in \(H\) containing \(a\), since planes
in the union of three cosets must take two points per coset.
Varying one vector puts it either outside \(H\), giving no
plane, or in \(E\) once, which gives directions avoiding \(a\).
A wedge with the last vector uses its component in \(E\) once,
giving directions in \(H\) avoiding \(a\). Varying the last
vector puts it in \(Q\), giving only directions containing \(a\).
Varying a leading vector gives two factors from \(Q^\perp\):
both in \(E\) give a direction containing \(a\), one outside
\(H\) gives no plane, and both outside \(H\) give a direction
outside \(H\). Thus the derivative and value have disjoint
direction supports in every case.
The uniform second derivative bound \(640h^2\) then gives
\(L_{0i}-2\le320h^2\). Using the weaker constant
\(h^2\le20\mathcal A\eta/g^2\) and each forward coefficient at
most \(2t_6\), the bracket is at most
\(25600t_6\mathcal A\eta/g^2\).

For \(\eta\ge\eta_0\), the universal \(L_{ij}\le5/2\) gives a
bracket at most \(2t_6\). Multiplying by \(168\) and subtracting
\(\eta\) in either case proves the comparison for
\[
 r_6(s)=\min\{s/2,\ s^4/40,\ \eta_0/336,\
                                      g^2/(4300800\mathcal A)\}.
\]
The paired reduction and stabilizer padding apply as above.
At a strict cutoff both marginal defects vanish at an optimum.
For distinct positive probabilities the same rigidity argument
gives precisely the local Clifford optimizers.
\end{proof}
